\documentclass[11pt]{article}
\usepackage[margin=1in]{geometry}
\usepackage{amsmath,amssymb,amsthm}
\usepackage{xurl}
\usepackage{hyperref}
\sloppy
\let\oldus\_ \renewcommand{\_}{\oldus\allowbreak}
\newtheorem{theorem}{Theorem}
\newtheorem{lemma}[theorem]{Lemma}
\theoremstyle{definition}
\newtheorem{definition}[theorem]{Definition}
\theoremstyle{remark}
\newtheorem{remark}[theorem]{Remark}
\newcommand{\OT}{\operatorname{OT}}
\newcommand{\rk}{r}
\newcommand{\depth}{\operatorname{depth}}

\title{Disproof of conjecture 00000000543:\\ for the three lines $x,y,x+y$ the Orlik--Terao algebra has
depth $\geq 1$, while $\mathrm{rank}-1-\beta=0$}
\author{Jackmeson1}
\date{2026-10-05}

\begin{document}
\maketitle

\section{The conjecture and the reading}

\textbf{Statement (English, verbatim).} Definition: Matroid complexes and their Orlik--Terao
algebras. Conjecture: The depth of the OT algebra equals the rank of the matroid minus 1 minus
$\beta$ (a depth-beta invariant formula for hyperplane arrangements). (The Chinese text states the
same formula: depth of the OT algebra $=$ rank of the matroid $-1-\beta$.)

\textbf{Reading.} For a (central) arrangement $A=\{H_1,\dots,H_m\}$ of distinct hyperplanes in
$k^n$, with matroid $M(A)$ on $E=\{1,\dots,m\}$, the claim is
\[
  (\ast)\qquad \depth \OT(A) \;=\; \rk(E) - 1 - \beta(M(A)).
\]
We work over an arbitrary field $k$ of characteristic $0$ (e.g.\ $\mathbb{Q}$). We give an
arrangement of three lines in $k^2$ for which $(\ast)$ fails with a non-negative right-hand side
(Theorem~\ref{thm:main}), and a second arrangement (four lines) for which the right-hand side is
$-1$, so that $(\ast)$ fails for \emph{every} notion of depth with values in
$\mathbb{N}\cup\{\infty\}$ (Theorem~\ref{thm:u24}).

\section{Definitions (as formalized)}

\begin{definition}[arrangement, matroid]
An arrangement of $m$ hyperplanes in $k^n$ is given by coefficient vectors
$v_1,\dots,v_m\in k^n$; the defining linear form of $H_i$ is
$\alpha_i=\sum_j v_i(j)\,x_j\in k[x_1,\dots,x_n]$. The rank function of the matroid is
$\rk(S)=\dim_k \operatorname{span}\{v_i : i\in S\}$ for $S\subseteq E$; the rank of the matroid is
$\rk(E)$. The beta invariant is
\[
  \beta \;=\; (-1)^{\rk(E)}\sum_{S\subseteq E}(-1)^{|S|}\,\rk(S).
\]
\end{definition}

\begin{definition}[Orlik--Terao algebra]
$\OT(A)$ is the $k$-subalgebra of the field of rational functions
$k(x_1,\dots,x_n)=\mathrm{Frac}\,k[x_1,\dots,x_n]$ generated by the reciprocals
$y_i=1/\alpha_i$. It is graded with every $y_i$ in degree one; since it is generated in degree one,
its irrelevant (homogeneous maximal) ideal is $\mathfrak{m}=(y_1,\dots,y_m)$.
\end{definition}

\begin{definition}[depth]
A list $(r_1,\dots,r_\ell)$ of elements of a ring $R$ is \emph{$R$-regular} if each $r_j$ is a
nonzerodivisor on $R/(r_1,\dots,r_{j-1})$ and $R/(r_1,\dots,r_\ell)\neq 0$ (Mathlib's
\texttt{RingTheory.Sequence.IsRegular}). We set
$\depth \OT(A)=\sup\{\ell : \text{there is an } \OT(A)\text{-regular sequence of length } \ell
\text{ in } \mathfrak{m}\}\in\mathbb{N}\cup\{\infty\}$. This is the same pattern as the depth
definition in our earlier package for conjecture 00000000537 (file \texttt{C537.lean}), which we
reuse with credit.
\end{definition}

\section{The counterexample}

Let $A$ be given by $v_1=(1,0)$, $v_2=(0,1)$, $v_3=(1,1)$, i.e.\ the lines $x=0$, $y=0$, $x+y=0$
in $k^2$. Its matroid is the uniform matroid $U_{2,3}$.

\begin{lemma}[rank and beta]\label{lem:rk}
If $v_1,\dots,v_m\in k^2$ are nonzero and pairwise linearly independent, then
$\rk(S)=\min(|S|,2)$ for all $S$. Hence for $A$: $\rk(E)=2$ and
$\beta=\sum_{j=0}^{3}\binom{3}{j}(-1)^j\min(j,2)=0-3+6-2=1$, so the right-hand side of
$(\ast)$ is $2-1-1=0$.
\end{lemma}

\begin{proof}
$\rk(S)\le\dim k^2=2$; if $|S|\ge 2$, $S$ contains two independent vectors, so $\rk(S)\ge 2$;
$\rk(\emptyset)=0$ and $\rk(\{i\})=1$ since $v_i\neq 0$. The sum depends only on $|S|$.
Pairwise independence of $(1,0),(0,1),(1,1)$ is checked coordinatewise.
\end{proof}

\begin{lemma}[$y_1$ is not a unit]\label{lem:unit}
Assume every $c_i:=\alpha_i(1,\dots,1)=\sum_j v_i(j)$ is nonzero (for $A$: $c=(1,1,2)$). Then no
$y_i$ is a unit of $\OT(A)$.
\end{lemma}

\begin{proof}
Let $Q=\prod_i\alpha_i\neq 0$ and $L=k[x][1/Q]\subseteq k(x)$ (Mathlib's
\texttt{Localization.subalgebra.ofField}). Since $y_i=(Q/\alpha_i)\cdot Q^{-1}$, every $y_i\in L$.
The $k$-algebra map $g:k[x]\to k[T,T^{-1}]$, $x_j\mapsto T^{-1}$, sends $\alpha_i\mapsto
c_iT^{-1}$, hence $Q$ to a unit, so it extends to $\mathrm{ev}:L\to k[T,T^{-1}]$, and
$\mathrm{ev}(y_i)=c_i^{-1}T$ (because $y_i\alpha_i=1$). By induction over the generators of
$\OT(A)$ (constants go to constants; sums and products are preserved), every $f\in\OT(A)$ lies
in $L$ and $\mathrm{ev}(f)$ lies in the image of $k[T]$. If $y_if=1$ with $f\in\OT(A)$, then
$c_i^{-1}T\,p(T)=1$ for some $p\in k[T]$; as $k[T]\to k[T,T^{-1}]$ is injective,
$T\cdot c_i^{-1}p=1$ in $k[T]$, so $T$ would be a unit of $k[T]$, which is false.
\end{proof}

\begin{theorem}\label{thm:main}
Under the hypothesis of Lemma~\ref{lem:unit}, $[y_i]$ is an $\OT(A)$-regular sequence in
$\mathfrak{m}$, so $\depth\OT(A)\ge 1$. For the arrangement $\{x,y,x+y\}$ this contradicts
$(\ast)$, whose right-hand side is $0$.
\end{theorem}

\begin{proof}
$\OT(A)$ is a subring of a field, hence a domain, and $y_i\neq0$ (as $\alpha_i\neq 0$), so $y_i$
is a nonzerodivisor. By Lemma~\ref{lem:unit}, $(y_i)\neq\OT(A)$, i.e.\ $\OT(A)/(y_i)\neq 0$. And
$y_i\in\mathfrak{m}$. So $\depth\OT(A)\ge 1>0=\rk(E)-1-\beta$ by Lemma~\ref{lem:rk}; in
particular $\depth\OT(A)$ is not equal to the integer $\rk(E)-1-\beta$.
\end{proof}

\begin{theorem}[a second witness]\label{thm:u24}
For the four lines $x,y,x+y,x+2y$ in $k^2$ (matroid $U_{2,4}$): $\rk(E)=2$,
$\beta=\sum_{j=0}^4\binom4j(-1)^j\min(j,2)=0-4+12-8+2=2$, and $\rk(E)-1-\beta=-1$. Hence no
$D\in\mathbb{N}\cup\{\infty\}$ equals the right-hand side of $(\ast)$, whatever reading of
``depth'' or of ``OT algebra'' is used.
\end{theorem}

\begin{proof}
Lemma~\ref{lem:rk} applies: the four vectors are nonzero and pairwise independent (in
characteristic $0$).
\end{proof}

\begin{remark}[context, not used in the proofs]
Proudfoot and Speyer, \emph{A broken circuit ring} (arXiv:math/0410069), study ``the ring
generated by the inverses of the restrictions of the linear functionals'' and state in their
abstract: ``In particular, R(L) is Cohen-Macaulay''. This is consistent with our finding that the
depth of $\OT(A)$ is positive; our proofs above are self-contained and do not rely on it.
\end{remark}

\section{Lean formalization}

File \texttt{lean/Conjecture543/Basic.lean} (namespace \texttt{C543}, only \texttt{import Mathlib}).
Definitions: \texttt{alpha} (written with the Greek letter in Lean), \texttt{y}, \texttt{OT} (\texttt{Algebra.adjoin k (Set.range (y v))} inside
\texttt{FractionRing (MvPolynomial (Fin n) k)}), \texttt{gen}, \texttt{irr}, \texttt{depth},
\texttt{rk}, \texttt{beta}; witnesses \texttt{vA} ($U_{2,3}$) and \texttt{vB} ($U_{2,4}$).
Main results:
\begin{itemize}
\item \texttt{not\_isUnit\_gen}, \texttt{isRegular\_gen}, \texttt{one\_le\_depth}: Lemma~\ref{lem:unit}
  and Theorem~\ref{thm:main} for any arrangement with all $c_i\neq0$.
\item \texttt{rk\_eq\_min}, \texttt{rk\_vA\_univ}, \texttt{beta\_vA}: Lemma~\ref{lem:rk}.
\item \texttt{witness}: $\rk(E)=2\wedge\beta=1\wedge 1\le\depth\wedge\neg\exists d\in\mathbb{N},\
  \depth=d\wedge d=\rk(E)-1-\beta$ for \texttt{vA}.
\item \texttt{conjecture\_543\_false}:
{\small
\begin{verbatim}
theorem conjecture_543_false :
  NOT forall (m n : Nat) (v : Fin m -> Fin n -> k),
    (forall i, v i != 0) ->
    (forall i j, i != j -> LinearIndependent k ![v i, v j]) ->
    rk v Finset.univ = n ->
    exists d : Nat, depth v = d /\
      (d : Int) = (rk v Finset.univ : Int) - 1 - beta v
\end{verbatim}
}
(ASCII transcription of the Lean statement.)
  (for every field $k$ of characteristic $0$; the hypotheses say: nonzero forms, distinct
  hyperplanes, essential arrangement).
\item \texttt{formula\_U24}, \texttt{not\_formula\_U24}: Theorem~\ref{thm:u24}.
\end{itemize}
\textbf{Axiom audit.} \texttt{\#print axioms} for \texttt{C543.conjecture\_543\_false},
\texttt{C543.witness} and \texttt{C543.not\_formula\_U24} reports only
\texttt{[propext, Classical.choice, Quot.sound]}. No \texttt{sorry}, no \texttt{native\_decide},
no added axioms.

\end{document}
